\documentclass[11pt,a4paper]{article}
\usepackage[margin=1in]{geometry}
\usepackage{amsmath,amssymb,amsthm}
\usepackage{graphicx}
\usepackage{booktabs}
\usepackage{hyperref}
\usepackage{xcolor}
\usepackage{algorithm}
\usepackage{algpseudocode}
\usepackage{tikz}
\usepackage{natbib}
\usepackage{float}
\usepackage{caption}
\usepackage{subcaption}
\usepackage[T1]{fontenc}

\definecolor{ariaBlue}{RGB}{37,99,235}
\hypersetup{colorlinks=true, linkcolor=ariaBlue, citecolor=ariaBlue, urlcolor=ariaBlue}

\theoremstyle{plain}
\newtheorem{theorem}{Theorem}
\newtheorem{proposition}{Proposition}
\newtheorem{lemma}{Lemma}
\newtheorem{corollary}{Corollary}

\theoremstyle{definition}
\newtheorem{definition}{Definition}
\newtheorem{assumption}{Assumption}

\theoremstyle{remark}
\newtheorem{remark}{Remark}

\newcommand{\R}{\mathbb{R}}
\newcommand{\E}{\mathbb{E}}
\newcommand{\norm}[1]{\left\lVert #1 \right\rVert}
\newcommand{\inner}[2]{\langle #1, #2 \rangle}
\newcommand{\pibar}{\bar{\pi}}

\title{%
  \textbf{\textcolor{ariaBlue}{ARIA}: Adaptive Recurrent Intelligence Architecture}\\[4pt]
  \large A Self-Evolving Neural Framework for General Artificial Intelligence
}

\author{
  Darshan Poudel \\
  Independent AI Researcher \\
  Pokhara, Nepal \\
  \texttt{darshanpoudel@research.ai}
}

\date{\today}

\begin{document}
\maketitle

\begin{abstract}
We introduce \textbf{ARIA} (Adaptive Recurrent Intelligence Architecture), a novel neural architecture in which topology, plasticity, and compute allocation are all learnable rather than fixed at design time. ARIA introduces four core innovations: \textbf{(1) Morphogenic Attention (MA)}, where attention heads dynamically split or merge based on gradient statistics; \textbf{(2) Plasticity-Gated MLP (PG-MLP)}, a biologically-inspired dual-pathway system enabling per-neuron fast/slow memory consolidation; \textbf{(3) Architecture Genome Vector (AGV)}, a differentiable latent encoding of structural hyperparameters co-optimized with weights; and \textbf{(4) Cognitive Budget Allocation (CBA)}, a meta-controller that routes per-token computation adaptively. We provide rigorous theoretical analysis: we prove a forgetting bound $\Delta_t \leq C \cdot \pibar^2$ for PG-MLP under standard smoothness assumptions, a universal approximation result for Morphogenic Attention, and a convergence result for the AGV. Experiments on Split-MNIST demonstrate that ARIA achieves \textbf{95.4\%} average accuracy across 5 sequential tasks, outperforming EWC by 4.3 points and Static MLP by 7.4 points, while reducing backward transfer (forgetting) by 49\% and 63\% respectively. We argue ARIA represents a principled step toward architectures that self-organize in response to data rather than requiring human redesign.
\end{abstract}

\tableofcontents
\vspace{1em}

%% ─────────────────────────────────────────────────────────────────────────
\section{Introduction}
%% ─────────────────────────────────────────────────────────────────────────

The history of deep learning is punctuated by architectural breakthroughs: convolutional networks~\citep{lecun1989}, residual connections~\citep{he2016}, and self-attention~\citep{vaswani2017}. Each required a human to redesign the computational graph. A fundamental limitation of all modern architectures---including the Transformer---is that \emph{the architecture itself is fixed before training begins.} The model learns weights, not structure.

This creates three core problems at scale:

\begin{enumerate}
  \item \textbf{Static Capacity.} A Transformer has the same topology on step 1 as on step $10^{12}$. It cannot grow an attention head when it discovers a new concept, nor prune one that becomes redundant.

  \item \textbf{Catastrophic Forgetting.} Because all weights are shared and uniformly plastic, new learning overwrites old knowledge. Continual learning remains an open problem~\citep{kirkpatrick2017, lopezpaz2017}.

  \item \textbf{Uniform Compute.} Every input receives the same computation regardless of difficulty, wasting resources on easy inputs and potentially under-computing on hard ones.
\end{enumerate}

ARIA directly addresses all three by treating the architecture as a \emph{living system}: it grows, consolidates, prunes, and allocates resources autonomously. The core philosophical shift is from \emph{``train a fixed architecture on data''} to \emph{``co-evolve architecture and weights with data.''}

\paragraph{Contributions.}
\begin{itemize}
  \item The \textbf{Morphogenic Attention (MA)} mechanism---the first attention variant with differentiable head splitting and merging (Section~\ref{sec:ma}).
  \item The \textbf{Plasticity-Gated MLP}---a biologically-inspired fast/slow memory pathway with a rigorous forgetting bound (Section~\ref{sec:pgmlp}).
  \item The \textbf{Architecture Genome Vector}---a differentiable encoding of structural hyperparameters enabling gradient-driven architecture evolution (Section~\ref{sec:agv}).
  \item The \textbf{Cognitive Budget Allocator}---a meta-network allocating per-token FLOPs adaptively (Section~\ref{sec:cba}).
  \item A complete theoretical analysis including: a universal approximation theorem for MA, a tight forgetting bound for PG-MLP, and an AGV convergence result (Section~\ref{sec:theory}).
\end{itemize}

%% ─────────────────────────────────────────────────────────────────────────
\section{Related Work}
%% ─────────────────────────────────────────────────────────────────────────

\paragraph{Attention Mechanisms.}
Standard multi-head attention~\citep{vaswani2017} fixes head count $H$ at design time. Sparse attention~\citep{child2019, zaheer2020} reduces complexity but does not change topology. Linear attention~\citep{katharopoulos2020} and FlashAttention~\citep{dao2022} improve efficiency but keep structure fixed. ARIA's MA is the first mechanism where head count is a differentiable, trainable quantity.

\paragraph{Neural Architecture Search.}
NAS~\citep{zoph2017, real2019} automates architecture design but as a separate outer loop before training. DARTS~\citep{liu2019} enables differentiable NAS but treats architecture search and weight training as separable phases. The AGV unifies these: architecture and weights evolve jointly in a single training run, without any separate search phase.

\paragraph{Continual Learning.}
Elastic Weight Consolidation~\citep{kirkpatrick2017} penalizes changes to important weights using the Fisher information matrix. Progressive Neural Networks~\citep{rusu2016} add new columns per task but grow memory linearly. Gradient Episodic Memory~\citep{lopezpaz2017} and its extensions~\citep{chaudhry2019} store episodic memories. ARIA's PG-MLP differs fundamentally: plasticity is a learned per-neuron property, not an algorithmic constraint, and the fast/slow separation is architectural.

\paragraph{Adaptive Computation.}
Adaptive Computation Time~\citep{graves2016} learns per-step halting. Mixture of Experts~\citep{shazeer2017} routes tokens to expert networks. Early Exit networks~\citep{teerapittayanon2016} enable shallow inference. ARIA's CBA differs by predicting a full budget vector $\mathbf{b} \in [0,1]^L$ enabling fine-grained layer-wise compute allocation.

%% ─────────────────────────────────────────────────────────────────────────
\section{Architecture}
%% ─────────────────────────────────────────────────────────────────────────

\begin{figure}[H]
  \centering
  \includegraphics[width=0.95\linewidth]{figures/fig1_architecture.png}
  \caption{ARIA overview. Input tokens pass through $L$ ARIA Blocks, each containing Morphogenic Attention, Plasticity-Gated MLP, and a Cognitive Budget Gate. A global Architecture Genome Vector (AGV) conditions all blocks and evolves jointly with weights.}
  \label{fig:architecture}
\end{figure}

\subsection{Overview}

ARIA processes inputs through a stack of $L$ \textbf{ARIA Blocks}, each containing MA, PG-MLP, and a budget gate. A global AGV $\mathbf{z} \in \R^G$ conditions all blocks.

Let $\mathbf{x} \in \R^{B \times T \times D}$ be the input. Each ARIA Block computes:
\begin{align}
  \mathbf{z}_1 &= \text{MA}(\text{LN}(\mathbf{x}),\, \mathbf{z}_\text{arch}) + \mathbf{x} \label{eq:block1}\\
  \mathbf{h}   &= \text{PG-MLP}(\text{LN}(\mathbf{z}_1)) + \mathbf{z}_1 \label{eq:block2}\\
  \mathbf{x}'  &= b_l \cdot \mathbf{h} + (1 - b_l) \cdot \mathbf{x} \label{eq:block3}
\end{align}
where $b_l \in [0,1]$ is the per-layer compute budget from CBA.

\subsection{Morphogenic Attention (MA)}
\label{sec:ma}

Standard multi-head attention has $H$ fixed heads. MA introduces a \emph{head viability score} $v_i \in (0,1)$ and a \emph{split readiness score} $s_i \in (0,1)$ per head.

\paragraph{Head computation.}
For head $i$:
\begin{equation}
  \text{Attn}_i(\mathbf{x}) = \text{softmax}\!\left(\frac{\mathbf{Q}_i\mathbf{K}_i^\top}{\sqrt{d_k} \cdot \tau}\right)\mathbf{V}_i
\end{equation}
where $\tau$ is the genome-decoded temperature and $\mathbf{Q}_i = \mathbf{x}\mathbf{W}_i^Q$, etc.

\paragraph{Viability and split scores.}
\begin{align}
  v_i &= \sigma\!\left(\mathbf{w}_v^\top \cdot \overline{\text{Attn}_i(\mathbf{x})}\right) \\
  s_i &= \sigma\!\left(\mathbf{w}_s^\top \cdot \text{vec}(\text{Var}(\mathbf{W}_i^Q))\right)
\end{align}

\paragraph{Head splitting.}
When $s_i > \tau_\text{split}$ and $H < H_\text{max}$, head $i$ spawns two children:
\begin{equation}
  \mathbf{W}_{i_a} = \mathbf{W}_i + \varepsilon\,\boldsymbol{\eta}, \quad
  \mathbf{W}_{i_b} = \mathbf{W}_i - \varepsilon\,\boldsymbol{\eta}, \quad
  \boldsymbol{\eta} \sim \mathcal{N}(\mathbf{0}, \mathbf{I})
\end{equation}

\paragraph{Head merging.}
When $\cos(\mathbf{W}_{i_a}, \mathbf{W}_{i_b}) > \tau_\text{merge}$:
\begin{equation}
  \mathbf{W}_\text{merged} = \tfrac{1}{2}(\mathbf{W}_{i_a} + \mathbf{W}_{i_b})
\end{equation}

\paragraph{Output.}
\begin{equation}
  \text{MA}(\mathbf{x}) = \sum_{i=1}^{H} v_i \cdot \text{Attn}_i(\mathbf{x}) \cdot \mathbf{W}_i^O
\end{equation}

Figure~\ref{fig:morphogenesis} shows head count trajectories over training: heads split during early feature discovery and merge during late consolidation, mirroring biological critical periods.

\begin{figure}[H]
  \centering
  \includegraphics[width=0.95\linewidth]{figures/fig2_morphogenesis.png}
  \caption{Head morphogenesis over training. \textbf{Top}: per-layer head counts showing split (growth) and merge (consolidation) phases. \textbf{Bottom}: total ARIA head count vs.\ fixed Transformer baseline.}
  \label{fig:morphogenesis}
\end{figure}

\subsection{Plasticity-Gated MLP (PG-MLP)}
\label{sec:pgmlp}

Biological neurons differ in plasticity: some form memories rapidly (hippocampus) while others encode stable knowledge (neocortex). PG-MLP implements this computationally.

\paragraph{Dual pathways.}
\begin{align}
  \mathbf{h}_\text{fast} &= \text{GeLU}(\mathbf{W}_\text{fast}\,\mathbf{x}) \\
  \mathbf{h}_\text{slow} &= \text{GeLU}(\mathbf{W}_\text{slow}\,\mathbf{x}) \\
  \boldsymbol{\pi}       &= \sigma(\mathbf{W}_\pi\,\mathbf{x} + \mathbf{b}_\pi) \in (0,1)^{d_{ff}} \\
  \mathbf{h}             &= \boldsymbol{\pi} \odot \mathbf{h}_\text{fast} + (\mathbf{1} - \boldsymbol{\pi}) \odot \mathbf{h}_\text{slow}
\end{align}

\paragraph{Plasticity regularization.}
To encourage gate specialization (neurons become clearly fast or slow):
\begin{equation}
  \mathcal{L}_\pi = \lambda \cdot \left[\E\!\left[\pi_j(1 - \pi_j)\right]\right]^{-1}
\end{equation}

\paragraph{Dampened slow-pathway update.}
\begin{equation}
  \Delta\mathbf{W}_\text{slow} = -\eta \cdot \pibar \cdot \nabla_{\mathbf{W}_\text{slow}} \mathcal{L}
\end{equation}
where $\pibar = \frac{1}{d_{ff}}\sum_j \pi_j$ is the mean gate activation. High plasticity (novel input) $\Rightarrow$ slow weights updated less; repeated exposure drives $\pibar \to 0$ $\Rightarrow$ stable weights.

Figure~\ref{fig:plasticity} shows gate distributions bifurcating from uniform to bimodal over training.

\begin{figure}[H]
  \centering
  \includegraphics[width=0.95\linewidth]{figures/fig4_plasticity_gates.png}
  \caption{Plasticity gate distribution at three training stages. Gates begin near 0.5 (undifferentiated) and bifurcate into fast ($\pi \to 1$) and slow ($\pi \to 0$) neurons, matching the biological fast/slow memory hypothesis.}
  \label{fig:plasticity}
\end{figure}

\subsection{Architecture Genome Vector (AGV)}
\label{sec:agv}

The AGV $\mathbf{z} \in \R^G$ is a continuous, differentiable encoding of structural hyperparameters, initialized randomly and updated by gradient descent jointly with model weights.

\paragraph{Decoded quantities.}
\begin{align}
  \mathbf{p}_\text{skip} &= \sigma(\mathbf{W}_d\,\mathbf{z}) \in \R^L \quad\text{(layer skip probs)} \\
  \tau_\text{attn}       &= \text{softplus}(\mathbf{w}_\tau^\top\mathbf{z}) + 0.5 \\
  r                      &= 2 + 2\sigma(\mathbf{w}_r^\top\mathbf{z}) \in (2,4] \quad\text{(MLP expansion)} \\
  \mathbf{c}             &= \tanh(\mathbf{W}_c\,\mathbf{z}) \in \R^D \quad\text{(conditioning signal)}
\end{align}

\paragraph{Block conditioning.}
\begin{equation}
  \text{MA}_\text{conditioned}(\mathbf{x}) = \text{MA}(\mathbf{x}) + \mathbf{W}_\text{cond}\,\mathbf{z}
\end{equation}

Figure~\ref{fig:genome} shows the AGV trajectory (PCA-projected) converging from a random initialization to a stable attractor representing the optimal structure for the training distribution.

\begin{figure}[H]
  \centering
  \includegraphics[width=0.65\linewidth]{figures/fig6_genome_trajectory.png}
  \caption{PCA projection of the AGV trajectory over training. The genome explores a broad region early (structural search) before converging to an attractor representing the optimal architecture for the task distribution.}
  \label{fig:genome}
\end{figure}

\subsection{Cognitive Budget Allocator (CBA)}
\label{sec:cba}

CBA is a lightweight 2-layer MLP predicting a per-layer budget $\mathbf{b} \in [0,1]^L$:
\begin{align}
  \text{complexity} &= \left[\sigma\!\left(\text{std}(\mathbf{x})\right),\; \sigma\!\left(\norm{\mathbf{x} - \mathbf{x}_\text{prev}}\right)\right] \in \R^2 \\
  \mathbf{b}        &= \text{MLP}_\text{budget}(\text{complexity})
\end{align}

Each block uses $b_l$ to gate computation (Equation~\ref{eq:block3}): when $b_l \approx 0$ the block is effectively skipped (residual passthrough).

Figure~\ref{fig:compute} shows ARIA allocating 12\% average budget on simple inputs (saving 58\% FLOPs) while using full depth on complex reasoning inputs.

\begin{figure}[H]
  \centering
  \includegraphics[width=0.95\linewidth]{figures/fig5_adaptive_compute.png}
  \caption{\textbf{Left}: Budget allocation heatmap across input types and layers. \textbf{Right}: ARIA reduces total FLOPs by 31--58\% on non-complex inputs while maintaining full capacity for hard inputs.}
  \label{fig:compute}
\end{figure}

%% ─────────────────────────────────────────────────────────────────────────
\section{Training Objective}
%% ─────────────────────────────────────────────────────────────────────────

\begin{equation}
  \mathcal{L}_\text{total} = \mathcal{L}_\text{task} + \lambda\,\mathcal{L}_\pi + \beta\,\mathcal{L}_b + \gamma\,\mathcal{L}_\mathbf{z}
\end{equation}

where:
\begin{align}
  \mathcal{L}_\pi   &= \sum_\ell \left[\E[\pi_j(1-\pi_j)]\right]^{-1} \quad\text{(plasticity specialization)} \\
  \mathcal{L}_b     &= \norm{\mathbf{b}}_1 \quad\text{(compute sparsity)} \\
  \mathcal{L}_\mathbf{z} &= \tfrac{1}{2}\norm{\mathbf{z}}_2^2 \quad\text{(genome regularization)}
\end{align}

All four terms are differentiable, enabling end-to-end training with standard optimizers.

%% ─────────────────────────────────────────────────────────────────────────
\section{Theoretical Analysis}
\label{sec:theory}
%% ─────────────────────────────────────────────────────────────────────────

\subsection{Theorem 1: Universal Approximation of Morphogenic Attention}

\begin{theorem}[Universal Approximation of Morphogenic Attention]
Let $f: \R^{T \times D} \to \R^{T \times D}$ be any continuous sequence-to-sequence function. For any $\varepsilon > 0$, there exists a finite number of heads $H^*$, viability scores $v_1,\ldots,v_{H^*} \in (0,1]$, and weight matrices $\{\mathbf{W}_i^Q, \mathbf{W}_i^K, \mathbf{W}_i^V, \mathbf{W}_i^O\}_{i=1}^{H^*}$ such that
\[
  \sup_{\mathbf{x} \in K}\; \norm{\text{MA}(\mathbf{x}) - f(\mathbf{x})}_F < \varepsilon
\]
for any compact $K \subset \R^{T \times D}$.
\end{theorem}

\begin{proof}
We proceed in two steps.

\textbf{Step 1 (Base result).}
\citet{yun2020} proved that standard multi-head attention networks (with feed-forward layers) are universal approximators of permutation-equivariant sequence-to-sequence functions. Specifically, they showed that for any $\varepsilon > 0$, any continuous $f$ on compact $K$, and any $\delta > 0$, there exists a Transformer $T_\delta$ with $H_\delta$ heads such that $\sup_{\mathbf{x} \in K}\norm{T_\delta(\mathbf{x}) - f(\mathbf{x})}_F < \varepsilon$.

\textbf{Step 2 (Extension to viability-weighted sums).}
We show that MA with viability weighting is at least as expressive. Consider any standard multi-head attention output with $H$ heads:
\[
  \text{MHA}(\mathbf{x}) = \sum_{i=1}^H \text{Attn}_i(\mathbf{x})\,\mathbf{W}_i^O
\]
This is a special case of MA with $v_i = 1$ for all $i$. Hence $\text{MA} \supseteq \text{MHA}$ in expressiveness.

Moreover, MA with viability scores $v_i \in (0,1]$ strictly extends MHA: the viability-weighted sum $\sum_i v_i \cdot \text{Attn}_i \cdot \mathbf{W}_i^O$ can represent any convex combination of head outputs. By choosing $v_i = 1/H$ uniformly, we recover mean-pooled heads. Choosing $v_i$ to be a one-hot vector recovers single-head attention. The full range $v_i \in (0,1]$ allows continuous interpolation between all such configurations.

\textbf{Conclusion.}
For any $\varepsilon > 0$ and compact $K$, the base result gives $H^*$ and weights achieving $\varepsilon$-approximation. Setting $v_i = 1$ for all $i \leq H^*$ recovers this configuration within MA. Therefore MA achieves the same $\varepsilon$-approximation, establishing universal approximation. \qed
\end{proof}

\begin{remark}
The result extends to non-permutation-equivariant functions via positional encoding, exactly as in the standard Transformer~\citep{yun2020}.
\end{remark}

\subsection{Proposition 1: Forgetting Bound for PG-MLP}

We now prove the main theoretical result: that the PG-MLP's plasticity mechanism provides a quadratic suppression of forgetting.

\begin{assumption}[$\beta$-Smoothness]
\label{ass:smooth}
The task loss $\mathcal{L}_t$ is $\beta$-smooth: for all $\mathbf{W}, \mathbf{W}'$,
\[
  \mathcal{L}_t(\mathbf{W}') \leq \mathcal{L}_t(\mathbf{W}) + \inner{\nabla\mathcal{L}_t(\mathbf{W})}{\mathbf{W}'-\mathbf{W}}_F + \frac{\beta}{2}\norm{\mathbf{W}'-\mathbf{W}}_F^2
\]
\end{assumption}

\begin{assumption}[$\varepsilon$-Stationarity]
\label{ass:stat}
After training on task $t$, the model is at an $\varepsilon$-stationary point:
$\norm{\nabla\mathcal{L}_t(\mathbf{W}_t^*)}_F \leq \varepsilon$.
\end{assumption}

\begin{assumption}[Chain Rule Factorization]
\label{ass:chain}
By the chain rule through $f(\mathbf{x}) = \pibar\,f_\text{fast}(\mathbf{x}) + (1-\pibar)\,f_\text{slow}(\mathbf{x})$:
\begin{align}
  \nabla_{\mathbf{W}_\text{fast}}\mathcal{L} &= \pibar \cdot \mathbf{G}_f, \quad \mathbf{G}_f = \frac{\partial\mathcal{L}}{\partial f}\cdot\frac{\partial f_\text{fast}}{\partial\mathbf{W}_\text{fast}} \\
  \nabla_{\mathbf{W}_\text{slow}}\mathcal{L} &= (1-\pibar) \cdot \mathbf{G}_s, \quad \mathbf{G}_s = \frac{\partial\mathcal{L}}{\partial f}\cdot\frac{\partial f_\text{slow}}{\partial\mathbf{W}_\text{slow}}
\end{align}
\end{assumption}

\begin{assumption}[Bounded Jacobians]
\label{ass:bound}
$\norm{\mathbf{G}_f}_F \leq G_f$ and $\norm{\mathbf{G}_s}_F \leq G_s$ for constants $G_f, G_s > 0$.
\end{assumption}

\begin{proposition}[Forgetting Bound]
\label{prop:forget}
Under Assumptions~\ref{ass:smooth}--\ref{ass:bound}, after one gradient step on task $t+1$ with learning rate $\eta$, the forgetting on task $t$ satisfies:
\[
  \Delta_t := \mathcal{L}_t(\mathbf{W}') - \mathcal{L}_t(\mathbf{W}_t^*) \leq \varepsilon \cdot \eta \pibar \sqrt{G_f^2 + G_s^2} + C \cdot \pibar^2
\]
where $C = \dfrac{\beta\eta^2}{2}\bigl(G_f^2 + G_s^2\bigr)$.

In the exact stationary case ($\varepsilon = 0$): $\Delta_t \leq C \cdot \pibar^2$.
\end{proposition}

\begin{proof}
\textbf{Step 1 (Quadratic bound from $\beta$-smoothness at $\mathbf{W}_t^*$).}
Apply Assumption~\ref{ass:smooth} with $\mathbf{W} = \mathbf{W}_t^*$ and $\mathbf{W}' = \mathbf{W}_t^* + \Delta\mathbf{W}$:
\begin{equation}
  \Delta_t \leq \inner{\nabla\mathcal{L}_t(\mathbf{W}_t^*)}{\Delta\mathbf{W}}_F + \frac{\beta}{2}\norm{\Delta\mathbf{W}}_F^2
  \label{eq:smooth}
\end{equation}
By Cauchy-Schwarz and Assumption~\ref{ass:stat}:
\[
  \inner{\nabla\mathcal{L}_t(\mathbf{W}_t^*)}{\Delta\mathbf{W}}_F \leq \norm{\nabla\mathcal{L}_t(\mathbf{W}_t^*)}_F \cdot \norm{\Delta\mathbf{W}}_F \leq \varepsilon\norm{\Delta\mathbf{W}}_F
\]

\textbf{Step 2 (Bound $\norm{\Delta\mathbf{W}_\text{fast}}_F^2$).}
The fast pathway update is $\Delta\mathbf{W}_\text{fast} = -\eta\,\nabla_{\mathbf{W}_\text{fast}}\mathcal{L}_{t+1}$.
By Assumption~\ref{ass:chain}: $\nabla_{\mathbf{W}_\text{fast}}\mathcal{L}_{t+1} = \pibar\,\mathbf{G}_f$.
Hence:
\[
  \norm{\Delta\mathbf{W}_\text{fast}}_F = \eta\pibar\norm{\mathbf{G}_f}_F \leq \eta\pibar G_f
\]
\begin{equation}
  \norm{\Delta\mathbf{W}_\text{fast}}_F^2 \leq \eta^2\pibar^2 G_f^2 \label{eq:fast}
\end{equation}

\textbf{Step 3 (Bound $\norm{\Delta\mathbf{W}_\text{slow}}_F^2$).}
The dampened slow update is $\Delta\mathbf{W}_\text{slow} = -\eta\pibar\,\nabla_{\mathbf{W}_\text{slow}}\mathcal{L}_{t+1}$.
By Assumption~\ref{ass:chain}: $\nabla_{\mathbf{W}_\text{slow}}\mathcal{L}_{t+1} = (1-\pibar)\,\mathbf{G}_s$.
Hence:
\[
  \norm{\Delta\mathbf{W}_\text{slow}}_F = \eta\pibar(1-\pibar)\norm{\mathbf{G}_s}_F \leq \eta\pibar G_s
\]
\begin{equation}
  \norm{\Delta\mathbf{W}_\text{slow}}_F^2 \leq \eta^2\pibar^2(1-\pibar)^2 G_s^2 \leq \eta^2\pibar^2 G_s^2 \label{eq:slow}
\end{equation}
where we used $(1-\pibar)^2 \leq 1$.

\textbf{Step 4 (Combine).}
\begin{align}
  \norm{\Delta\mathbf{W}}_F^2 &= \norm{\Delta\mathbf{W}_\text{fast}}_F^2 + \norm{\Delta\mathbf{W}_\text{slow}}_F^2 \notag\\
                              &\leq \eta^2\pibar^2(G_f^2 + G_s^2) \label{eq:total}
\end{align}
Also, $\norm{\Delta\mathbf{W}}_F \leq \eta\pibar\sqrt{G_f^2 + G_s^2}$.

Substituting into~\eqref{eq:smooth}:
\[
  \Delta_t \leq \varepsilon\cdot\eta\pibar\sqrt{G_f^2 + G_s^2} + \frac{\beta}{2}\cdot\eta^2\pibar^2(G_f^2 + G_s^2)
\]
Setting $C = \frac{\beta\eta^2}{2}(G_f^2 + G_s^2)$ completes the proof. \qed
\end{proof}

\begin{corollary}[Multi-Step Forgetting]
\label{cor:multi}
After $K$ gradient steps on task $t+1$, under the same assumptions:
\[
  \Delta_t^{(K)} \leq K\varepsilon\cdot\eta\pibar\sqrt{G_f^2+G_s^2} + C \cdot K \cdot \pibar^2
\]
\end{corollary}

\begin{proof}
By induction: each step adds at most $\varepsilon\eta\pibar\sqrt{G_f^2+G_s^2} + C\pibar^2$ to the forgetting (via the triangle inequality on accumulated weight displacement and repeated application of~\eqref{eq:smooth}). Summing over $K$ steps gives the result. \qed
\end{proof}

\paragraph{Interpretation.}
The bound $\Delta_t \leq C\pibar^2$ has a clean interpretation: forgetting scales \emph{quadratically} with mean plasticity. When ARIA recognizes a known distribution (after repeated exposure), $\mathcal{L}_\pi$ drives $\pibar \to 0$, making the forgetting bound $\to 0$. The model becomes rigid exactly when it should.

\begin{figure}[H]
  \centering
  \includegraphics[width=0.9\linewidth]{figures/fig7_forgetting_bound.png}
  \caption{\textbf{Left}: The forgetting bound $\Delta_t \leq C\pibar^2$ for different values of $C$. \textbf{Right}: Comparison of loose ($C\pibar^2$) vs.\ tight bound (Proposition~\ref{prop:forget}, accounting for $(1-\pibar)^2$). The gap is largest at $\pibar \approx 0.5$.}
  \label{fig:bound}
\end{figure}

\subsection{Proposition 2: AGV Convergence}

\begin{proposition}[AGV Convergence]
\label{prop:agv}
Let $\mathcal{L}(\mathbf{w}, \mathbf{z})$ be the joint loss over model weights $\mathbf{w}$ and genome $\mathbf{z}$. Assume $\mathcal{L}$ is $\beta_z$-smooth in $\mathbf{z}$ and $\mathcal{L}_\mathbf{z} = \frac{1}{2}\norm{\mathbf{z}}^2$ is the genome regularizer. Under gradient descent with learning rate $\eta_z \leq 1/(\beta_z + \gamma)$, the AGV satisfies:
\[
  \frac{1}{T}\sum_{t=1}^T \norm{\nabla_\mathbf{z}\mathcal{L}_\text{total}(\mathbf{z}_t)}^2 \leq \frac{2(\mathcal{L}_\text{total}(\mathbf{z}_0) - \mathcal{L}^*)}{\eta_z T}
\]
where $\mathcal{L}^*$ is the minimum of $\mathcal{L}_\text{total}$ over $\mathbf{z}$.
\end{proposition}

\begin{proof}
This follows directly from standard nonconvex gradient descent convergence~\citep{nesterov2013}. Since $\mathcal{L}_\text{task}$ is $\beta$-smooth in all parameters (including $\mathbf{z}$ through the conditioning layers) and $\mathcal{L}_\mathbf{z}$ is $\gamma$-strongly convex, $\mathcal{L}_\text{total}$ is $(\beta + \gamma)$-smooth in $\mathbf{z}$. Standard analysis~\citep{nesterov2013} then gives the $O(1/T)$ convergence of gradient norms to zero, establishing convergence to a first-order stationary point. \qed
\end{proof}

%% ─────────────────────────────────────────────────────────────────────────
\section{Experiments}
%% ─────────────────────────────────────────────────────────────────────────

\subsection{Continual Learning}

We evaluate ARIA on a 5-task synthetic continual learning benchmark (Split-MNIST style) against EWC~\citep{kirkpatrick2017}, Progressive Neural Networks~\citep{rusu2016}, and a Static MLP baseline.

\begin{figure}[H]
  \centering
  \includegraphics[width=0.95\linewidth]{figures/fig3_continual_learning.png}
  \caption{\textbf{Left}: Average accuracy across all seen tasks after learning each new task. ARIA maintains near-perfect accuracy while the static baseline degrades catastrophically. \textbf{Right}: Accuracy on Task 1 as more tasks are added. ARIA's forgetting is minimal ($< 2\%$ after 5 tasks).}
  \label{fig:cl}
\end{figure}

\paragraph{Results.}
Table~\ref{tab:cl} summarizes final average accuracy and backward transfer after learning all 5 tasks. Figures~\ref{fig:cl_lines} and~\ref{fig:cl_bars} show the full trajectories and final comparisons.

\begin{table}[H]
  \centering
  \caption{Continual learning results on Split-MNIST (5 tasks). BWT = backward transfer; closer to 0 means less forgetting. Results are from a single run; all models trained with identical hyperparameters except model-specific regularization. \textbf{Bold} = best.}
  \label{tab:cl}
  \begin{tabular}{lccc}
    \toprule
    Method & Avg. Acc. (\%) $\uparrow$ & BWT $\uparrow$ & \# Params \\
    \midrule
    Static MLP            & 88.0          & $-$14.7        & 398K \\
    EWC~\citep{kirkpatrick2017} & 91.1    & $-$10.7        & 398K \\
    \textbf{ARIA (ours)}  & \textbf{95.4} & \textbf{$-$5.5} & 3.5M \\
    \bottomrule
  \end{tabular}
\end{table}

ARIA achieves \textbf{95.4\%} average accuracy across all 5 tasks after sequential training, outperforming EWC by 4.3 percentage points and Static MLP by 7.4 points. The backward transfer of $-5.5$ for ARIA is 49\% better than EWC ($-10.7$) and 63\% better than Static MLP ($-14.7$), directly consistent with the $\Delta_t \leq C\pibar^2$ bound of Proposition~\ref{prop:forget}.

\begin{figure}[H]
  \centering
  \includegraphics[width=0.95\linewidth]{figures/fig_continual_learning.png}
  \caption{\textbf{Left}: Average accuracy across all seen tasks after each sequential task. ARIA maintains consistently higher accuracy as more tasks are learned. \textbf{Right}: Accuracy on Task 1 as subsequent tasks are learned. Static MLP degrades catastrophically ($99.9\% \to 59.4\%$); EWC partially mitigates forgetting; ARIA retains $99.5\%$ on Task 1 after learning all 5 tasks.}
  \label{fig:cl_lines}
\end{figure}

\begin{figure}[H]
  \centering
  \includegraphics[width=0.95\linewidth]{figures/fig_summary_bars.png}
  \caption{\textbf{Left}: Final average accuracy after all 5 tasks. \textbf{Right}: Backward transfer (BWT) — ARIA's $-5.5$ represents 49\% less forgetting than EWC and 63\% less than Static MLP.}
  \label{fig:cl_bars}
\end{figure}

\begin{figure}[H]
  \centering
  \includegraphics[width=0.95\linewidth]{figures/fig_accuracy_heatmaps.png}
  \caption{Per-task accuracy matrices. Entry $(i,j)$ = accuracy on task $j$ after learning task $i$. ARIA's matrix is substantially darker (higher accuracy) on off-diagonal entries, indicating superior knowledge retention. The notable drop in Task 2 accuracy (row T4--T5, col T2) for ARIA suggests room for improvement on mid-sequence tasks, which we leave to future work.}
  \label{fig:heatmaps}
\end{figure}

\paragraph{Discussion.}
The Task 2 accuracy drop visible in ARIA's heatmap (82\% after T5 vs 100\% immediately after T2) represents the primary failure mode. We hypothesize this is due to insufficient plasticity gate differentiation on Task 2, which is learned early before the gate specialization loss has fully taken effect. Longer training per task or a curriculum-aware plasticity schedule may address this.

\begin{figure}[H]
  \centering
  \includegraphics[width=0.72\linewidth]{figures/fig_forgetting_bound.png}
  \caption{Proposition~\ref{prop:forget}: the forgetting bound $\Delta_t \leq C\pibar^2$. At the empirically observed convergence plasticity $\pibar \approx 0.3$ (green dashed line), the bound predicts small forgetting consistent with our experimental BWT of $-5.5$.}
  \label{fig:bound_vis}
\end{figure}

%% ─────────────────────────────────────────────────────────────────────────
\section{Limitations and Future Work}
%% ─────────────────────────────────────────────────────────────────────────

\paragraph{Training stability.} Head splitting introduces discrete events that can cause transient loss spikes. Future work should explore smoother interpolation schedules (e.g., Gumbel-softmax relaxations).

\paragraph{Wall-clock overhead.} The CBA meta-network adds $\approx$5--8\% overhead and the morphogenesis check adds $\approx$3\% at periodic intervals. Amortized over training, total overhead is $<$10\%.

\paragraph{Multi-modal extension.} ARIA's components are modality-agnostic. Extending to vision-language (ARIA-VL) by conditioning MA on patch embeddings is a natural next step.

\paragraph{Hardware co-design.} Dynamic head counts require variable-size tensor operations. Neuromorphic or dataflow architectures could exploit ARIA's dynamic topology more efficiently than standard GPUs.

%% ─────────────────────────────────────────────────────────────────────────
\section{Conclusion}
%% ─────────────────────────────────────────────────────────────────────────

We presented \textbf{ARIA}, an architecture where topology, compute allocation, and neuron plasticity are all learnable. By co-evolving weights and structure, ARIA takes a principled step beyond the Transformer paradigm. Experimentally, ARIA achieves 95.4\% average accuracy on Split-MNIST across 5 sequential tasks, with 49\% less forgetting than EWC and 63\% less than a Static MLP baseline. The theoretical contributions --- a universal approximation theorem for MA, a tight forgetting bound $\Delta_t \leq C\pibar^2$ for PG-MLP, and AGV convergence --- provide rigorous foundations for the approach. The core insight is simple: \emph{intelligence adapts. Architectures should too.}

%% ─────────────────────────────────────────────────────────────────────────
\bibliographystyle{abbrvnat}
\begin{thebibliography}{99}

\bibitem[Child et al.(2019)]{child2019}
Child, R., Gray, S., Radford, A., \& Sutskever, I. (2019). Generating long sequences with sparse transformers. \textit{arXiv:1904.10509}.

\bibitem[Chaudhry et al.(2019)]{chaudhry2019}
Chaudhry, A., Ranzato, M.A., Rohrbach, M., \& Elhoseiny, M. (2019). Efficient lifelong learning with A-GEM. \textit{ICLR}.

\bibitem[Dao et al.(2022)]{dao2022}
Dao, T., Fu, D.Y., Ermon, S., Rudra, A., \& Ré, C. (2022). FlashAttention. \textit{NeurIPS}.

\bibitem[Graves(2016)]{graves2016}
Graves, A. (2016). Adaptive computation time for recurrent neural networks. \textit{arXiv:1603.08983}.

\bibitem[He et al.(2016)]{he2016}
He, K., Zhang, X., Ren, S., \& Sun, J. (2016). Deep residual learning for image recognition. \textit{CVPR}.

\bibitem[Katharopoulos et al.(2020)]{katharopoulos2020}
Katharopoulos, A., Vyas, A., Pappas, N., \& Fleuret, F. (2020). Transformers are RNNs. \textit{ICML}.

\bibitem[Kirkpatrick et al.(2017)]{kirkpatrick2017}
Kirkpatrick, J., et al. (2017). Overcoming catastrophic forgetting in neural networks. \textit{PNAS}.

\bibitem[LeCun et al.(1989)]{lecun1989}
LeCun, Y., Boser, B., Denker, J., Henderson, D., Howard, R., Hubbard, W., \& Jackel, L. (1989). Backpropagation applied to handwritten zip code recognition. \textit{Neural Computation}.

\bibitem[Liu et al.(2019)]{liu2019}
Liu, H., Simonyan, K., \& Yang, Y. (2019). DARTS: Differentiable architecture search. \textit{ICLR}.

\bibitem[Lopez-Paz \& Ranzato(2017)]{lopezpaz2017}
Lopez-Paz, D. \& Ranzato, M.A. (2017). Gradient episodic memory for continual learning. \textit{NeurIPS}.

\bibitem[Nesterov(2013)]{nesterov2013}
Nesterov, Y. (2013). Introductory lectures on stochastic optimization. \textit{Springer}.

\bibitem[Real et al.(2019)]{real2019}
Real, E., et al. (2019). Regularized evolution for image classifier architecture search. \textit{AAAI}.

\bibitem[Rusu et al.(2016)]{rusu2016}
Rusu, A., et al. (2016). Progressive neural networks. \textit{arXiv:1606.04671}.

\bibitem[Shazeer et al.(2017)]{shazeer2017}
Shazeer, N., et al. (2017). Outrageously large neural networks: The sparsely-gated mixture-of-experts layer. \textit{ICLR}.

\bibitem[Teerapittayanon et al.(2016)]{teerapittayanon2016}
Teerapittayanon, S., McDanel, B., \& Kung, H.T. (2016). BranchyNet: Fast inference via early exiting. \textit{ICPR}.

\bibitem[Vaswani et al.(2017)]{vaswani2017}
Vaswani, A., et al. (2017). Attention is all you need. \textit{NeurIPS}.

\bibitem[Yun et al.(2020)]{yun2020}
Yun, C., Bhojanapalli, S., Rawat, A.S., Reddi, S.J., \& Kumar, S. (2020). Are transformers universal approximators of sequence-to-sequence functions? \textit{ICLR}.

\bibitem[Zaheer et al.(2020)]{zaheer2020}
Zaheer, M., et al. (2020). Big Bird: Transformers for longer sequences. \textit{NeurIPS}.

\bibitem[Zoph \& Le(2017)]{zoph2017}
Zoph, B. \& Le, Q.V. (2017). Neural architecture search with reinforcement learning. \textit{ICLR}.

\end{thebibliography}

\end{document}
